\documentclass[12pt]{article}
\usepackage[utf8]{inputenc}
\usepackage[T1]{fontenc}
\usepackage[margin=1in]{geometry}
\usepackage{setspace}
\usepackage{xcolor}
\usepackage[colorlinks=true, urlcolor=blue!60!black, linkcolor=blue!60!black]{hyperref}
\setlength{\parindent}{2em}
\setlength{\parskip}{1.25em}
\renewcommand{\baselinestretch}{1.0}
\begin{document}
\begin{center}
{\LARGE \textit{Ex parte McCardle}, 74 U.S. (7 Wall.) 506 (1869)\par}
\vspace{0.8em}
{\large Case Brief --- Four-Page Edition\par}
\end{center}
\vspace{1em}

\section{Memory Jogger}

Congress may withdraw a statutory route of Supreme Court appellate jurisdiction even after a case has been argued: a unanimous Court dismissed a Reconstruction-era habeas appeal once Congress repealed the 1867 appellate provision, while noting that its older habeas jurisdiction under the Judiciary Act of 1789 remained intact (\href{https://supreme.justia.com/cases/federal/us/74/506/}{74 U.S.\ 506}).

\section{Detailed Case Facts}

The Military Reconstruction Acts of 1867 placed the former Confederate states under military government. William H.\ McCardle, a civilian newspaper editor in Vicksburg, Mississippi, published editorials sharply denouncing Reconstruction. Federal military authorities arrested him and held him for trial before a military commission on charges arising from allegedly incendiary and libelous articles.

McCardle petitioned the federal circuit court for habeas corpus under the Habeas Corpus Act of February 5, 1867, which authorized the writ for anyone restrained in violation of the Constitution, treaties, or laws of the United States and expressly allowed appeals to the Supreme Court.

\section{Procedural History}

The circuit court remanded McCardle to military custody, and he appealed under the 1867 Act. The Supreme Court denied the government's motion to dismiss (73 U.S.\ 318 (1868)) and heard four days of argument on the merits in March 1868. Fearing the Court would invalidate the Reconstruction Acts, as \textit{Ex parte Milligan} suggested it might, Congress on March 27, 1868, over President Johnson's veto, repealed the 1867 Act insofar as it authorized appeals to the Supreme Court, including appeals already taken. The Court postponed decision, over a protest by Justices Grier and Field, while Chief Justice Chase presided over Johnson's impeachment trial, then heard reargument on the repeal's effect.

\section{Judicial Votes}

Dismissed for want of jurisdiction, \textbf{8--0}, on April 12, 1869.
\begin{itemize}
\item \textbf{Opinion of the Court:} Chief Justice Chase, joined by Justices Nelson, Grier, Clifford, Swayne, Miller, Davis, and Field.
\item \textbf{Concurring or dissenting:} None.
\end{itemize}

\section{Holding}

Under the Exceptions Clause of Article III, Section 2, Congress could repeal the 1867 grant of appellate jurisdiction even while McCardle's appeal was pending after argument, and the Court may not inquire into Congress's motives. Once statutory jurisdiction is repealed, the Court cannot pronounce judgment and must dismiss. The 1868 Act removed only appeals under the 1867 Act and left intact the Court's preexisting habeas jurisdiction under the Judiciary Act of 1789.

\section{Analysis of Opinions}

\textbf{Majority (Chase).} The opinion proceeds in four steps. \emph{Source of jurisdiction:} the Court's appellate jurisdiction comes from Article III, not statutes, but it is conferred ``with such Exceptions, and under such Regulations as the Congress shall make.'' Under \textit{Durousseau} and \textit{Wiscart}, affirmative grants imply exceptions for what is not granted; here the exception was express, ``hardly possible to imagine a plainer instance of positive exception.'' \emph{Motive:} the Court examines only power, and the power to make exceptions is granted in express words. \emph{Pending appeals:} unlike legislative interference with a court's continuing jurisdiction, repeal of a jurisdictional statute means it is treated, except as to closed transactions, as if it never existed (\textit{Norris v.\ Crocker}; \textit{Insurance Co.\ v.\ Ritchie}). \emph{Alternative jurisdiction:} counsel wrongly assumed that the Court lost all habeas appellate power; the 1868 Act did ``not affect the jurisdiction which was previously exercised'' under the 1789 Act.

\textbf{Subsequent developments.} Months later, \textit{Ex parte Yerger} (1869) confirmed that the 1789 route remained open, and \textit{Felker v.\ Turpin} (1996) relied on \textit{Yerger} to uphold AEDPA's bar on certiorari review of gatekeeping orders. \textit{United States v.\ Klein} (1871) held that Congress may not use a jurisdictional strip to prescribe a ``rule of decision.'' \textit{Boumediene v.\ Bush} (2008) held that Congress cannot strip habeas jurisdiction without an adequate substitute. A plurality in \textit{Patchak v.\ Zinke} (2018) reaffirmed \textit{McCardle} over Chief Justice Roberts's \textit{Klein}-based dissent.

\section{Examples: Future Applications}

\begin{enumerate}
\item \textbf{Same-Side (Upheld): Repealing a Direct-Appeal Route.} Congress repeals direct appeals to the Supreme Court in a class of cases, including pending ones, but leaves appeal to a court of appeals and certiorari. Pending direct appeals must be dismissed, and another door remains open.
\item \textbf{Same-Side (Upheld): Barring Interlocutory Certiorari.} Barring certiorari on interlocutory class-certification appeals while preserving review after final judgment regulates timing, as in \textit{Felker}.
\item \textbf{Opposite-Side (Likely Invalid): Outcome-Selective Strip.} A statute barring all Supreme Court review, including original habeas, only of decisions \emph{upholding} flag-burning convictions eliminates the \textit{Yerger} escape valve and rigs outcomes in violation of \textit{Klein}, the First and Fifth Amendments, and the Suspension Clause.
\item \textbf{Opposite-Side (Invalid): Detainee Habeas Without Substitute.} Stripping habeas for detainees at a base under complete U.S.\ control, with only an inadequate status review, fails under \textit{Boumediene}.
\item \textbf{Fence-Sitter: Stripping All Review of School-Prayer Cases.} A facially neutral bar on any Supreme Court review of school-prayer challenges, leaving lower courts open, fits \textit{McCardle}'s text but leaves no alternative route and squarely presents Henry Hart's ``essential functions'' objection. The Court has not resolved the question.
\end{enumerate}

\section{Critique}

\textbf{Broad versus narrow reading.} Read broadly, \textit{McCardle}'s statement that the power to make exceptions ``is given by express words'' suggests plenary congressional power to remove Supreme Court review of controversial constitutional subjects such as school prayer, abortion, or gun regulation. Read narrowly, it holds only that Congress may close one statutory avenue while another, here the 1789 Act, remains open.

\textbf{The ``essential functions'' thesis.} Henry Hart argued that exceptions cannot ``engulf the rule'' or destroy the Court's essential role in keeping federal law supreme and uniform. Critics respond that Article III's text contains no such carve-out and that the Madisonian Compromise left even the lower federal courts to congressional discretion.

\section{Key Quotations}

\begin{quote}
``Without jurisdiction, the court cannot proceed at all in any cause. Jurisdiction is power to declare the law, and, when it ceases to exist, the only function remaining to the court is that of announcing the fact and dismissing the cause.''
\end{quote}
\noindent --- Chief Justice Chase, Opinion of the Court, 74 U.S.\ at 514.

\begin{quote}
``We are not at liberty to inquire into the motives of the legislature. We can only examine into its power under the Constitution, and the power to make exceptions to the appellate jurisdiction of this court is given by express words.''
\end{quote}
\noindent --- Chief Justice Chase, Opinion of the Court, 74 U.S.\ at 514.

\begin{quote}
``Counsel seem to have supposed, if effect be given to the repealing act in question, that the whole appellate power of the court, in cases of habeas corpus, is denied. But this is an error. The act of 1868 does not except from that jurisdiction any cases but appeals from Circuit Courts under the act of 1867. It does not affect the jurisdiction which was previously exercised.''
\end{quote}
\noindent --- Chief Justice Chase, Opinion of the Court, 74 U.S.\ at 515.

\end{document}
